%In this section, we first introduce the target microservice system in Alibaba, and then explain the availability issues, anomaly types, and metrics considered in this paper.

%\subsection{Target System}
The e-commerce system of Alibaba has more than 846 millions monthly active users.
It adopts the microservice architecture and contains more than 30,000 services.
These microservices are deployed with containers (i.e., Docker~\cite{docker}) and virtual machines and orchestrated by the customized orchestrator based on Kubernetes~\cite{kubernetes}, with Service Mesh~\cite{istio} applied for service communications.
The system is equipped with a large-scale monitoring infrastructure called EagleEye~\cite{EagleEye}.
EagleEye includes a tracing SDK, a real-time ETL (Extract-Transform-Load) data processing system, a batch processing computing cluster, and a web-based user interface.

%However, the enormous traffic, complicated business process and frequent invocation between microserivices make the system fragile and hard to locate the root causes of failures, so there is a great challenge for diagnosing the failures in large-scale systems efficiently and accurately.
As the carrier of the business, the system needs to ensure high availability.
Therefore, a business monitoring platform is deployed to raise timely alarms about availability issues.
These availability issues usually indicate problems with the running of business, for example the drop of successfully placed orders and success rate of transactions.
An availability issue can be caused by different types of anomalies, each of which is indicated by a set of metrics.
An anomaly can originate from a service and propagate along service calls, and finally cause an availability issue.
In this work, we focus on the following three types of anomalies that cause most of the availability issues in Alibaba.

\begin{itemize}
\item  \textbf{Performance Anomaly}.
Performance anomaly is indicated by anomalous increase of response time (RT).
It is usually caused by problematic implementation or improper environmental configurations (e.g., CPU/memory configurations of containers and virtual machines).

\item  \textbf{Reliability Anomaly}.
Reliability anomaly is indicated by anomalous increase of error counts (EC), i.e., the numbers of service call failures.
It is usually caused by exceptions due to code defects or environmental failures (e.g., server or network outage).

\item  \textbf{Traffic Anomaly}.
Traffic anomaly is indicated by anomalous increase or decrease of queries per second (QPS).
Anomalous traffic increase may break the services, while anomalous decrease may indicate that many requests cannot reach the services.
Traffic anomaly is usually caused by improper traffic configurations (e.g., the traffic limits of Nginx~\cite{nginx}), DoS attack, or unanticipated stress test.

\end{itemize}

In anomaly detection, the 3-sigma rule is often used to identify outliers of metric values as candidates of anomalies. 
The 3-sigma rule states that in a normal distribution almost all the values remain within three standard deviations of the mean.
The range can be represented as ($\mu- 3\sigma$, $\mu + 3\sigma$), where $\mu$ is the mean and $\sigma$ is the standard deviation.
The values within the range account for 99.73\% of all the values and the others can be regarded as outliers.

%The prerequisite of \app~mainly contain three parts:
%1) Service calls, which are used to construct the service call graph.
%2) Metrics, which are used to enrich the service call graph, detect and rank the root cause microservices.
%3) Single root cause, supposing that an available issue is caused by only one root cause.
